% !TEX root = ../main.tex
\section{De 5 incitamentkonflikter mellom eiere og virksomhetsledere}\label{thr:incentive-conflicts-owner-manager}

\paragraph{Definisjon.} BSZ kap.~10 (slides) lister fem klassiske \emph{incentive conflicts} som oppstår fordi eierne (principalen) og virksomhetslederen (agenten) har forskjellige mål, ulik informasjon og ulik tidshorisont og risikoprofil. Konfliktene er kjernen i hvorfor virksomheten -- som et \emph{nexus of contracts} -- må designe kontrakter, måltall og insentiver. [SLIDE]

\subsection*{Forklaring -- de fem konfliktene}

\subsubsection*{1. Work hard or shirk (innsats / shirking)}
Agentens innsats skaper verdi for eieren, men er en kostnad for agenten selv (anstrengelse, tapt fritid). Når innsats ikke kan observeres perfekt, vil agenten typisk velge lavere innsats enn det som er optimalt for eieren. Dette er kjernen i \thref{moral-hazard}{moral hazard}: hidden action etter kontraktinngåelsen. [SLIDE]

\emph{Eksempel:} Lederen bruker arbeidstiden på å spille golf med kunder framfor faglig krevende analyser, og scenen er vanskelig å observere ovenfra.

\subsubsection*{2. Salaries and perks (lønn og personlige goder)}
Agenten ønsker høy lønn \emph{og} ikke-pekuniære goder -- firmajet, prestisjekontorer, kunst på veggene, konferanser i eksotiske destinasjoner. Disse perks koster ofte virksomheten mer enn de er verdt for agenten (en flytur til \$50\,000 gir kanskje agenten \$10\,000 i opplevd verdi), og er derfor en ren overføring fra eier til leder med dødvektstap. [SLIDE]

\emph{Eksempel:} Tycos Dennis Kozlowski-skandalen (2002) med firmamidler brukt på private fester og kunstinnkjøp.

\subsubsection*{3. Risk aversion (risikoaversjon)}
Lederen har mye av sin formue og humankapital konsentrert i ett selskap (job, bonus, opsjoner, omdømme). Eierne kan derimot diversifisere sine porteføljer. Lederen er derfor mer risikoavers enn eieren og kan velge bort prosjekter med positiv NPV men høy varians, eller velge for konservative finansieringsstrukturer (lav gjeldsgrad). Eier vil at lederen tar på seg \emph{all} risiko med positiv forventet verdi. [SLIDE]

\emph{Eksempel:} En CEO unngår å satse på ny teknologi fordi en mislykket satsing kan gi sparken, selv om forventet avkastning er attraktiv for diversifiserte aksjonærer.

\subsubsection*{4. Differential time horizons (horisontproblemet)}
Lederen har endelig kontraktsperiode (typisk 4--6 år som CEO) og pensjonering i sikte; eierne (særlig institusjonelle) har i prinsippet evig horisont via aksjen. Lederen kan derfor velge prosjekter med tidlige cash-flows og avvise prosjekter med stor verdi langt fram i tid. Bonus knyttet til årlig regnskapsmessig profit forsterker problemet. [SLIDE]

\emph{Eksempel:} En CEO kutter R\&D og vedlikehold de siste 2 årene før pensjon for å booste kortsiktig EBIT og bonus, vel vitende om at problemene rammer etterfølgeren.

\subsubsection*{5. Overinvestment / empire building}
Større virksomhet $\Rightarrow$ høyere status, høyere lønn, flere underordnede, mer ``free cash'' å disponere. Lederen kan derfor være fristet til oppkjøp og ekspansjon \emph{utover} det som maksimerer aksjonærverdi -- klassisk \emph{empire building} (Jensen 1986: free cash flow-problemet). Speilsiden er at lederen kan være motvillig til å \emph{krympe} virksomheten selv når avvikling er optimalt. [SLIDE/BOK]

\emph{Eksempel:} Verdi\-ødelggende oppkjøp drevet av status og bonus knyttet til omsetning eller selskapsstørrelse.

\subsection*{Hvor opptrer konfliktene?}

Slidene understreker at konfliktene ikke bare oppstår mellom eier$\to$leder, men også:
\begin{itemize}
\item \textbf{Leder $\to$ medarbeider} -- shirking, perks i miniatyr, kortsiktig handlingshorisont innen avdelingen.
\item \textbf{Virksomhet $\to$ leverandør} -- leverandøren shirker på kvalitet, leverer billige innsatsfaktorer.
\item \textbf{Aksjonær $\to$ kreditor} -- ``asset substitution'' (eierne velger risikable prosjekter på kreditors regning).
\end{itemize}
Konfliktene rammer relasjonene i ulik grad og krever derfor ulike kontraktsdesign. [SLIDE]

\subsection*{Kjerneforutsetninger}
\begin{itemize}
\item Asymmetrisk informasjon mellom eier og leder (innsats, evner, prosjektkvalitet).
\item Agenten er risikoavers; eieren er diversifisert.
\item Agenten har endelig horisont; eieren har lengre/uendelig horisont.
\item Lederens nytte øker med både lønn, perks og selskapsstørrelse.
\item Kontrakter er ufullstendige -- ikke alle kontingenser kan beskrives.
\end{itemize}

\subsection*{Muligheter (hva forklaringen brukes til)}
\begin{itemize}
\item Begrunner aksjebasert avlønning (kobler lederens formue til aksjekurs $\Rightarrow$ adresserer pkt.\ 3, 4).
\item Begrunner LTI (long-term incentives) og vesting over flere år (adresserer pkt.\ 4).
\item Begrunner monitoring via styret, eksterne revisorer, analytikere (adresserer pkt.\ 1, 2, 5).
\item Begrunner gjeld som disiplinerende mekanisme (Jensen: gjeld disiplinerer free cash flow $\Rightarrow$ pkt.\ 5).
\item Forklarer hvorfor eierkonsentrasjon (familieselskaper, PE-fond) ofte korrelerer med færre agency-problemer.
\end{itemize}

\subsection*{Begrensninger og kritikk}
\begin{itemize}
\item Listen er ikke uttømmende -- f.eks.\ \emph{differential treatment} (favorisering av enkelte underordnede) og \emph{asset substitution} mot kreditorer er også relevante.
\item Konfliktene overlapper og forsterker hverandre (risikoaversjon + horisont $\Rightarrow$ underinvestering).
\item Empirisk er sammenhengen mellom CEO-insentiver og firm performance svakere enn teorien forutsier.
\item Sterk fokus på finansielle insentiver kan fortrenge intrinsisk motivasjon og samarbeidskultur.
\item Listen er ``vestlig'' børsnotert--orientert; i familieselskaper, koop\-erativer eller statlige selskaper kan andre konflikter dominere.
\end{itemize}

\subsection*{Modell / oversikt}

\begin{center}
\renewcommand{\arraystretch}{1.15}
\begin{tabular}{@{}p{3.4cm}p{4.0cm}p{4.4cm}p{3.4cm}@{}}
\toprule
\textbf{Konflikt} & \textbf{Hva agenten gjør} & \textbf{Kilde} & \textbf{Vanlig kontrakts\-løsning} \\
\midrule
1. Shirking & For lav innsats & Hidden action & Insentivlønn, monitoring \\
2. Perks & Overforbruk av frynsegoder & Eier betaler, agent nyter & Tak på frynsegoder, eier-overvåking \\
3. Risikoaversjon & Avviser god risikabel investering & Udiversifisert humankapital & Aksjeopsjoner (konvekse insentiver) \\
4. Horisont & Kortsiktig profit, kutt R\&D & Kort kontraktsperiode & LTI med lang vesting, klausuler ved fratreden \\
5. Empire building & Verdi\-ødelggende vekst & Status- og lønnskoblet til størrelse & Gjeld, ROIC-mål, M\&A-disiplin \\
\bottomrule
\end{tabular}
\end{center}

\subsection*{Eksempler}
\begin{itemize}
\item \textbf{Equinor (Statoil) -- Hydro fusjon (2007):} Sammenslåingen av Statoils olje\-virksomhet med Norsk Hydros petroleums\-divisjon ble bredt fortolket som verdi-skapende industriell logikk, men kritikere pekte på \emph{empire building}-elementer i den nye konsernsjefens domene. [EKS]
\item \textbf{Vy/NSB -- bonusbråk (2019):} Konsernsjef Geir Isaksens bonus etter omorganisering ble offentlig kontroversiell -- klassisk konflikt om hvordan agenten belønnes for verdi som faktisk er drevet av eksogene faktorer (kontraktstildelinger fra staten). [EKS]
\item \textbf{Volkswagen Dieselgate (2015):} Kombinasjon av risiko (juks for å nå mål), horisontproblem (kortsiktig salg) og shirking på etisk standard. Eierne (porsche/piech-familien + tyske institusjoner) bar tap på \euro 30 mrd.+. [EKS]
\item \textbf{Norsk Hydro / Alunorte (2018):} Lokale agenter under\-rapporterte miljørisiko ved Alunorte-raffineriet i Brasil -- shirking på rapporteringsplikt, og selskapet betalte stor reputasjons- og finansiell pris. [EKS]
\end{itemize}

\subsection*{Kobling til søyle og andre teorier}
Søyle: \SOne\ (eier delegerer beslutningsrett -- konflikten oppstår i delegasjonen), \STwo\ (måltallene som skal disiplinere agenten), \SThree\ (insentivlønnen som skal lukke gapet). Relaterte teorier: \thref{principal-agent}{Principal-agent-teori}, \thref{moral-hazard}{Moral hazard} (pkt.\ 1), \thref{adverse-selection}{Adverse selection} (informasjons\-problemet bak pkt.\ 3), \thref{costless-contracting}{Costless contracting} (benchmark uten konfliktene), \thref{costly-contracting}{Costly contracting} (realistisk verden), \thref{agency-costs}{Agency costs} (kostnaden ved konfliktene), \thref{game-theory}{Spilteori} (analyse av monitor/shirk-spill og repeterte interaksjoner).

\paragraph{Brukt i spørsmål:} Q03, Q04.
